

\section{From Counterfactuals To Treatment Effects}\label{sec:effect}

We finally turn our attention to intervals for ITE for subjects not in
the study, and for which both potential oucomes are missing.

\subsection{A naive approach}\label{subsec:naive}

Consider an arbitrary testing point $x$.  Having constructed tools for
counterfactual inference producing interval estimates of a given
potential outcome, we can construct a pair of prediction intervals at
level $1-\alpha/2$, namely, $[\hat{Y}^{L}(1; x), \hat{Y}^{R}(1; x)]$
for $Y(1)$ and $[\hat{Y}^{L}(0; x), \hat{Y}^{R}(0; x)]$ for $Y(0)$. By
contrasting these two intervals, we can obtain an interval for ITE as
follows: 
\[\CITE(x) = [\hat{Y}^{L}(1; x) - \hat{Y}^{R}(0; x), \hat{Y}^{R}(1; x) - \hat{Y}^{L}(0; x)].\]
If the counterfactual intervals have guaranteed coverage, then
$\CITE(x)$ also covers ITE at the level $1-\alpha$ in the sense of
\eqref{eq:covITE}. 
% \[
% \P_{X\sim P_{X}}(Y(1) - Y(0)\in \hat{C}(X))\ge 1 - \alpha.
% \]
This means that we can use
counterfactual intervals produced by weighted split-CQR or by BART
whenever they are suitable.

\subsection{A nested approach}\label{subsec:nest}
% Despite its simplicity, the naive approach
% entirely decouples two
% potential outcomes and thus 
% can be conservative.
In this subsection,
we propose another strategy which we refer to as the nested approach. The nested procedure
starts by splitting the data into two folds. On the first fold we
train $\hat{C}_{1}(x)$ and $\hat{C}_{0}(x)$ by applying counterfactual
inference. On the second fold, for each unit $i$, we compute
$\hat{C}_{0}(X_{i})$ if $T_{i} = 1$ and compute $\hat{C}_{1}(X_{i})$
if $T_{i} = 0$. This induces an interval $\CITE(x; t, y^{\obs})$ for
ITE defined as
\[
\CITE(x; t, y^{\obs}) = 
    \begin{cases}
      y^{\obs} - \hat{C}_{0}(x), & t = 1,\\
      \hat{C}_{1}(x) - y^{\obs}, & t = 0. 
    \end{cases}
 \]
 Put $\hat{C}_{i} = \CITE(X_{i}; T_{i}, Y_{i}^{\obs})$  as in Table
 \ref{tab:twofolds}, which illustrates the procedure. For any unit $i$
 in the second fold,
\[
 \P(Y_{i}(1) - Y_{i}(0)\in \hat{C}_{i}) 
 = \P(T_{i} = 1)\P(Y_{i}(0)\in \hat{C}_{0}(X_{i}) \mid T_{i} = 1) + \P(T_{i} = 0)\P(Y_{i}(1)\in \hat{C}_{1}(X_{i}) \mid T_{i} = 0).
\]
This gives that if
\begin{equation}
  \label{eq:condition}
  \P(Y_{i}(0)\in \hat{C}_{0}(X_{i}) \mid T_{i} = 1) \ge 1 - \alpha, \quad \P(Y_{i}(1)\in \hat{C}_{1}(X_{i}) \mid T_{i} = 0)\ge 1 - \alpha,
\end{equation}
then
\begin{equation}
  \label{eq:fold2_cond}
  \P(Y_{i}(1) - Y_{i}(0)\in \hat{C}_{i})\ge 1 - \alpha.
\end{equation}


\begin{table}
  \caption{\label{tab:twofolds} Sketch of the two folds in the nested approach: the
    first fold (left) is used to construct counterfactual intervals;
    the second fold (right) is used to infer ITE.}  

  \scriptsize
  \begin{tabular}{ccc}
    \toprule
    $Y_{i}(1)$ & $Y_{i}(0)$ & $Y_{i}^{\mathrm{obs}}$\\
    \midrule
    \multicolumn{3}{l}{Treatment Group (fold 1)}\\
    \cmark & \xmark & $Y_{i}(1)$\\
    \midrule
    \multicolumn{3}{l}{Control Group (fold 1)}\\
    \xmark & \cmark & $Y_{i}(0)$\\
    \bottomrule
  \end{tabular}
  \hspace{0.5cm}
  \scriptsize  
  \begin{tabular}{ccc}
    \toprule
    $\hat{Y}_{i}(1)$ & $\hat{Y}_{i}(0)$ & $\hat{C}_{i}$ \\
    \midrule
    \multicolumn{3}{l}{Treatment Group (fold 2)}\\
    $Y_{i}(1)$ & $[\hat{Y}^{L}_{i}(0), \hat{Y}^{R}_{i}(0)]$ & $[Y_{i}(1) - \hat{Y}^{R}_{i}(0), Y_{i}(1) - \hat{Y}^{L}_{i}(0)]$\\
    \midrule
    \multicolumn{3}{l}{Control Group (fold 2)}\\
    $[\hat{Y}^{L}_{i}(1), \hat{Y}^{R}_{i}(1)]$ & $Y_{i}(0)$ & $[\hat{Y}^{L}_{i}(1) - Y_{i}(0), \hat{Y}^{R}_{i}(1) - Y_{i}(0)]$\\
    \bottomrule
  \end{tabular}
\end{table}


For randomized experiments with known propensity score $e(x)$,
\eqref{eq:condition} is satisfied if we use
$w_{0}(x) = e(x) / (1 - e(x))$ for $\hat{C}_{0}(x)$ and use
$w_{1}(x) = (1 - e(x)) / e(x)$ for $\hat{C}_{1}(x)$. Please note that
we do not need to split $\alpha$ for $\hat{C}_{0}(x)$ and
$\hat{C}_{1}(x)$. For observational studies, we can substitute $e(x)$
with $\hat{e}(x)$. By Theorem \ref{thm:double_robustness_ATE},
\eqref{eq:condition} holds approximately if either the propensity
scores or the conditional quantiles are estimated well.

% In general, using the same argument as above, for any given covariate distribution $Q$ which may differ from $P_{X}$,
% \[\P_{X\sim Q}(Y(1) - Y(0) \in \hat{C}(X; T, Y^{\obs}))\ge 1 - \alpha\]
% if the following conditions on $\hat{C}_{1}(x)$ and $\hat{C}_{0}(x)$ holds. 
% \[\P_{X\sim Q}(Y(0)\in \hat{C}_{0}(X) \mid T = 1) \ge 1 - \alpha, \quad \P_{X\sim Q}(Y(1)\in \hat{C}_{1}(X) \mid T = 0)\ge 1 - \alpha.\]
% By Table \ref{tab:weight_func}, this can be achieved by setting $w_{0}(x) =  / (1 - e(x))$ for $\hat{C}_{0}(x)$ and use $w_{1}(x) = (1 - e(x)) / e(x)$ for $\hat{C}_{1}(x)$

The nested procedure creates an i.i.d.~dataset $(X_{i}, \hat{C}_{i})$,
conditional on the first fold, such that
$\hat{C}_{i} = \CITE(X_{i}; T_{i}, Y_{i}^{\obs})$ covers the ITE
with probability at least $1 - \alpha$. Therefore, the intervals
$\hat{C}_{i}$ serve as ``surrogate intervals'' for ITE. If we can fit
a model of $\hat{C}_{i}$ on $X_{i}$, denoted by $\tdCITE(x)$, then the
intervals can be generalized to subjects with both potential outcomes
missing. If $\tdCITE(X_{i})$ does not shrink $\hat{C}_{i}$ drastically,
it is likely that $\P(Y(1) - Y(0)\in \tdCITE(X_{i}))$ shall be close to
or above $1 - \alpha$.

It may be useful to think of the nested method as follows:
  instead of estimating the unobserved uncertainty, the nested method
  is fitting an observed ``uncertainty measurement''
  $\hat{C}_{i}$. This is arguably simpler.

\subsection{An inexact and an exact method under the nested framework}
The function $\tdCITE(x)$ can be obtained by training a model of the
left- and right-end point of $\hat{C}_{i}$ on $X_{i}$ separately using
generic machine learning methods. For instance, we can fit the 40\% quantile of the left-end point and 60\% quantile of the right-end point. Note that $\tdCITE$ does not have the
same theoretical guarantee as those offered by conformalized
counterfactual inference since the fit may not be controlled. For this
reason, we refer to it as an ``inexact method''. Nonetheless,  as will
be shown later, the inexact method achieves, in our empirical
examples, the target coverage with drastically shorter intervals than
the naive approach.


In extremely sensitive settings where the validity of predictions is
of serious concern, we may still need a method with a theoretical
guarantee of coverage. Here, we propose a secondary conformal
inference procedure on the induced dataset $(X_{i}, \hat{C}_{i})$.
Given a generic observation $(X, T, Y^{\obs})$, let
$C = \CITE(X, T, Y^{\obs})$ denote the induced interval. The second
procedure---the ``exact method''---is based on the simple observation
that if one can find an interval expansion function $\hat{\cC}(\cdot)$
that maps a covariate value to an interval such that
\begin{equation}
  \label{eq:interval_conformal_goal}
  \P(C\subset \hat{\cC}_\text{ITE}(X))\ge 1 - \gamma,
\end{equation}
then by \eqref{eq:fold2_cond}, 
\[\P(Y(1) - Y(0) \not\in \hat{\cC}_\text{ITE}(X))\le \P(Y(1) - Y(0) \not\in C) + \P(C\not\subset \hat{\cC}_\text{ITE}(X))\le \alpha + \gamma.\]

\begin{algorithm}[H]
  \caption{(Unweighted) conformal inference for interval outcomes}  \label{algo:int}
  \textbf{Input: } level $\gamma$, data $\Z = (X_{i}, C_{i})_{i\in \I}$ where $C_{i} = [C_{i}^{L}, C_{i}^{R}]$, testing point $x$, 

  \hspace{1.3cm} functions $\hat{m}^{L}(x; \D), \hat{m}^{R}(x; \D)$ to fit the conditional mean/median of $C^{L}, C^{R}$

  \textbf{Procedure: }
  \begin{algorithmic}[1]
    \State Split $\Z$ into a training fold $\Z_{\tr} \triangleq (X_{i}, C_{i})_{i\in \I_{\tr}}$ and a calibration fold $\Z_{\ca}\triangleq (X_{i}, C_{i})_{i\in \I_{\ca}}$
    \State For each $i \in \I_{\ca}$, compute score $V_{i} = \max\{\hat{m}^{L}(X_{i}; \Z_{\tr}) - C_{i}^{L}, C_{i}^{R} - \hat{m}^{R}(X_{i}; \Z_{\tr})\}$
    \State Compute $\eta$ as the $(1 - \gamma)(1 + 1 / |\Z_{\ca}|)$ quantile of the empirical distribution of $\{V_{i}: i\in \I_{\ca}\}$
  \end{algorithmic}

  \textbf{Output: $\hat{\cC}(x) = [\hat{m}^{L}(x; \Z_{\tr}) - \eta, \hat{m}^{R}(x; \Z_{\tr}) + \eta]$}
\end{algorithm}


Denote by $C^{L}$ and $C^{R}$ the left- and the right-end point of
$C$, respectively. To achieve \eqref{eq:interval_conformal_goal}, we
need to find a lower prediction bound for $C^{L}$ and an upper
prediction bound for $C^{R}$. A naive approach is to apply standard
(unweighted) one-sided conformal inference on $C^{L}$ and $C^{R}$ with
level $\gamma / 2$ separately. Such a crude Bonferroni correction may
be conservative in practice. To overcome this, we propose a conformal
inference procedure that jointly calibrates $C^{L}$ and $C^{R}$. The
coverage guarantee can be proved using the standard argument
\citep[e.g.][]{lei2018distribution, romano2019conformalized} and we
include the proof in Appendix \ref{app:miscellaneous} for
completeness. 
\begin{theorem}\label{thm:int_conformal}
  Consider Algorithm \ref{algo:int} and assume $(X_{i}, C_{i})\stackrel{i.i.d.}{\sim} (X, C)$. Then
  \[\P(C\subset \hat{\cC}(X))\ge 1 - \gamma.\]
\end{theorem}

With all this in place, both the inexact and the exact methods are stated in
Algorithm \ref{algo:two_step}.
\begin{algorithm}
  \caption{Nested approach for interval estimates of ITE}  \label{algo:two_step}
  \textbf{Input: } level $\alpha$, level $\gamma$ (only for exact version), data $\Z = (X_{i}, Y_{i}, T_{i})_{i=1}^{n}$, testing point $x$

  \vspace{0.2cm}
  \textbf{Step I: data splitting}
  \begin{algorithmic}[1]
    \State Split the data into two folds $\Z_{1}$ and $\Z_{2}$
    \State Estimate propensity score $\hat{e}(x)$ on $\Z_{1}$
  \end{algorithmic}

  \vspace{0.2cm}
  \textbf{Step II: counterfactual inference on $\Z_{2}$}
  \begin{algorithmic}[1]   
    \For{$i$ in $\Z_{2}$ with $T_{i} = 1$}
         \State Compute $[\hat{Y}_{i}^{L}(0), \hat{Y}_{i}^{R}(0)]$ in Algorithm \ref{algo:split-CQR} on $\Z_{1}$ with level $\alpha$ and $w_{0}(x) = \hat{e}(x) / (1 - \hat{e}(x))$
         \State Compute $\hat{C}_{i} = [Y_{i}(1) - \hat{Y}_{i}^{R}(0), Y_{i}(1) - \hat{Y}_{i}^{L}(0)]$
    \EndFor
    \For{$i$ in $\Z_{2}$ with $T_{i} = 0$}
         \State Compute $[\hat{Y}_{i}^{L}(1), \hat{Y}_{i}^{R}(1)]$ in Algorithm \ref{algo:split-CQR} on $\Z_{1}$ with level $\alpha$ and $w_{1}(x) = (1 - \hat{e}(x)) / \hat{e}(x)$
         \State Compute $\hat{C}_{i} = [\hat{Y}_{i}^{L}(1) - Y_{i}(0), \hat{Y}_{i}^{R}(1) - Y_{i}(0)]$
    \EndFor
  \end{algorithmic}

  \vspace{0.2cm}
  \textbf{Step III: Interval of ITE on the testing point}
  \begin{algorithmic}[1]
    \State (Exact version) Apply Algorithm \ref{algo:int} on $(X_{i}, \hat{C}_{i})_{i\in \Z_{2}}$ with level $\gamma$, yielding an interval $\hat{\cC}_\text{ITE}(x)$
    % \State (Inexact version) Fit the conditional mean/median $\hat{m}^{L}(x)$ of $\hat{C}^{L}$ and $\hat{m}^{R}(x)$ of $\hat{C}^{R}$ separately and construct $\hat{\cC}_\text{ITE}(x) = [\hat{m}^{L}(x), \hat{m}^{R}(x)]$
    \State (Inexact version) Fit conditional quantiles of $\hat{C}^{L}$ and $\hat{C}^{R}$, yielding an interval $\hat{\cC}_\text{ITE}(x)$
  \end{algorithmic}
  \textbf{Output: $\hat{\cC}_\text{ITE}(x)$}\label{algo:nest_CQR}
\end{algorithm}


\subsection{Empirical performance}

To evaluate the performance of our methods, we design numerical
experiments on the data analyzed in the 2018 Atlantic Causal Inference
Conference workshop on heterogeneous treatment effects
\citep{carvalho2019assessing}. The workshop organizers generated a
synthetic dataset based on the National Study of Learning Mindsets
(NLSM) \citep{yeager2019national}, a large-scale randomized trial of a
behavioral intervention, to emulate an observational study.  For
information on the dataset, please see Section 2 from
\cite{carvalho2019assessing}. Due to privacy concerns, the workshop
organizers only released limited information on the data generating
process, as well as the simulated dataset, available at
\url{https://github.com/grf-labs/grf/tree/master/experiments/acic18}.
Although the focus of this workshop was
on heterogeneous treatment effects, the organizers did not evaluate
whether the submissions cover the ITE or CATE. In this subsection, we
shall fill in this gap by comparing our method with Causal Forest,
X-learner and BART, just as we did in Section \ref{subsec:simul}.

Obviously, we must know the ground truth in order to evaluate
coverage. Therefore, we generated synthetic datasets based on the
available information from \cite{carvalho2019assessing}. First, we
split data into two folds $\Z_{1}$ and $\Z_{2}$, with
$|\Z_{1}| = 2079$ including 20\% of the samples and $|\Z_{2}| = 8312$
including the remaining 80\%. In our numerical experiments, we
generate the covariate vector $X$ by sampling from $\Z_{2}$ with
replacement. To generate the potential outcomes, we apply random
forest from \texttt{R} \texttt{randomForest} package on $\Z_{1}$ to
fit $\E [Y(0)]$. Denote the output by $\hat{m}_{0}(x)$. Then we
generate $\E [Y(1)]$ by adding the CATE function $\tau(x)$ (equation
(1) of \cite{carvalho2019assessing}) to $\hat{m}_{0}(x)$. To account
for heteroscedasticity, we apply quantile random forest from
\texttt{R} \texttt{grf} package to fit the 25\% and 75\% conditional
quantiles of $Y(0)$ and $Y(1)$ and compute the conditional
interquartile ranges $\hat{r}_{0}(x)$ and $\hat{r}_{1}(x)$. Given a
covariate vector $X_{i}$, we subsequently generate $Y_{i}(1)$ and
$Y_{i}(0)$ as
\[
Y_{i}(1) = \hat{m}_{0}(X_{i}) + \tau(X_{i}) +
0.5\hat{r}_{1}(X_{i})\eps_{i1}, \quad Y_{i}(0) = \hat{m}_{0}(X_{i}) +
0.5\hat{r}_{0}(X_{i})\eps_{i0},\quad \eps_{i1},
\eps_{i0}\stackrel{i.i.d.}{\sim} N(0, 1).
\]
Finally, we generate propensity scores by applying random forest from
\texttt{R} \texttt{randomForest} package on $\Z_{1}$ and truncating
the estimated propensity score $\hat{e}(x)$ at $0.1$ and $0.9$ to
guarantee overlap. For each $X_{i}$, we generate $T_{i}$ as a
Bernoulli random variable with parameter $\hat{e}(X_{i})$.

For each run, we first generate $6000$ quadruples
$(X_{i}, T_{i}, Y_{i}(1), Y_{i}(0))$ sampled according to the above
data generating process. We then randomly select $n = 1000$ samples as
the training set to induce an observational study with observations
$(X_{i}, T_{i}, Y_{i}^{\obs})$. For the remaining $5000$ testing
samples, only the covariates $X_{i}$ are accessible to the analyst;
the triple $(T_{i}, Y_{i}(1), Y_{i}(0))$ is solely used to evaluate
coverage. It goes without saying that each method will take the
training set as the input and, for each testing sample, produce $95\%$
intervals for ITE.

For weighted split-CQR, we apply the naive procedure as described in Section \ref{subsec:naive} and the nested procedure with both exact and inexact calibration as described in Algorithm \ref{algo:nest_CQR}. For all procedures, we apply BART to fit conditional quantiles, as in Section \ref{subsec:simul}, and apply gradient boosting from \texttt{R} \texttt{gbm} package to fit propensity scores. For the exact nested method, we set $\alpha = \gamma = 0.025$. % All methods are implemented in \texttt{R} \texttt{cfcausal} package, available at \url{https://github.com/lihualei71/cfcausal}.

Note that neither the naive nor the nested approaches are limited to
CQR. Therefore, we also wrap both methods around BART as competitors,
where only the inexact version is considered for the nested approach
since BART cannot produce exact counterfactual intervals. Finally,
since Causal Forest and X-learner cannot produce counterfactual
intervals, except in the special case from Section \ref{subsec:simul}
where $Y(0)\equiv 0$, we do not wrap the naive or the nested methods
around them but directly report their confidence intervals as
benchmarks instead. % The R programs to replicate the results are available at \url{https://github.com/lihualei71/cfcausalPaper}.

\begin{figure}[H]
  \centering
  \includegraphics[width = 0.48\textwidth]{{simul_NLSM_coverage_paper}.pdf}
  \includegraphics[width = 0.48\textwidth]{{simul_NLSM_len_paper}.pdf}
  \caption{Coverage (left) and average length (right) of intervals for ITE on synthetic data generated from the NLSM data. The red vertical line corresponds to the target coverage $95\%$ and the blue vertical line corresponds to the average length of oracle intervals.}\label{fig:NLSM_simul_ITE}
\end{figure}

\begin{figure}[H]
  \centering
  \includegraphics[width = 0.9\textwidth]{simul_NLSM_cond_paper.pdf}
  \caption{Estimated conditional coverage of ITE as a function of the
    conditional variance $\sigma^{2}(x)$ (upper row), and CATE
    $\tau(x)$ (lower row). The coverage is set to $95\%$. The blue
    curves correspond to the median and the blue confidence bands are
    the $5\%$ and $95\%$ quantiles of these estimates across $100$
    replicates.}\label{fig:NLSM_cond}

\end{figure}


Figure \ref{fig:NLSM_simul_ITE} presents the coverage and the average
length of intervals estimated on the testing set by repeating the
above procedure 100 times. As expected, the naive methods with both
CQR and BART are conservative. The exact nested method with CQR is
also conservative although the coverage is only guaranteed to be above
$95\%$ in the worst case. In contrast, the inexact nested methods with
CQR or BART are less conservative. Notably, BART fails to achieve the
desired coverage but CQR, with BART as the learner, calibrates it
successfully. Also, we can see from the right panel that the average
length of intervals of inexact-CQR is just slightly above that of
inexact-BART while significantly lower than that of either the naive
or the exact nested methods. Moreover, in accordance with the
observations from Section \ref{subsec:simul}, Causal Forest and
X-learner have poor coverage and, thus, their (short) intervals are
misleading. As in Section \ref{subsec:simul}, we compute the average
length of oracle intervals as $3.92\E [\sigma(X)]$ where
$\sigma(x) = 0.5\sqrt{\hat{r}_{1}(x)^{2} + \hat{r}_{0}(x)^{2}}$.
Having said this, we are in a different situation here since the true
conditional quantiles of ITE can never be identified without
assumptions on the joint distribution of $Y(1)$ and $Y(0)$.
Therefore, this oracle length cannot be achieved in general no matter
how powerful the model fitting. Keeping this cautionary remark in
mind, we nonetheless see that weighted split-CQR still produces
reasonably short intervals.

Finally, as in Section \ref{subsec:simul}, we investigate the
conditional coverage as a function of the conditional variance
$\sigma^{2}(x)$ and the CATE $\tau(x)$, respectively. For better
visualization, we exclude Causal Forest and X-learner since they have
poor marginal coverage. We can see from Figure \ref{fig:NLSM_cond}
that inexact-CQR has desirable and relatively even conditional
coverage while inexact-BART performs worse.


\subsection{Re-analysing NLSM data}
In this subsection, we apply inexact-CQR with BART as the learner to re-analyze the NLSM data from \cite{carvalho2019assessing}. Since the ground truth is inaccessible, we only perform an exploratory analysis for the purpose of illustration.

To create informative testing points, we split the data into two folds
$\Z_{1}$ and $\Z_{2}$. Then we apply weighted inexact-CQR on $\Z_{1}$
to produce intervals for ITE for each point in $\Z_{2}$. Similarly, we
apply weighted inexact-CQR on $\Z_{2}$ to produce intervals of ITE for
each point in $\Z_{1}$. To account for the variability from data
splitting, we repeat the above procedure $100$ times. Figure
\ref{fig:NLSM_analysis} (a) displays the average length of intervals
as a function of level $\alpha$ with the upper and lower envelopes
being respectively the $95\%$ and $5\%$ quantiles across $100$ runs.

\begin{figure}
  \centering
  \includegraphics[width = 0.32\textwidth]{analysis_NLSM_len_paper.pdf}
  \includegraphics[width = 0.32\textwidth]{analysis_NLSM_pos_paper.pdf}
  \includegraphics[width = 0.32\textwidth]{analysis_NLSM_neg_paper.pdf}
  \caption{Re-analysis of NLSM data: (a) average length of intervals; (b) fraction of intervals with positive lower bounds; (c) fraction of intervals with negative upper bounds}\label{fig:NLSM_analysis}  
\end{figure}


With these intervals, we can make inferential claims on ITE with
confidence. For instance, we can decide to assign treatment to a
patient if the lower prediction bound of her ITE interval is
positive. Since these intervals are guaranteed to have desired
coverage, it is likely that they produce fewer false positives on
  the average, although we leave a theoretical investigation of this
intuition to future work. Figure \ref{fig:NLSM_analysis} (b) and (c)
show the fractions of intervals that only cover positive and negative
values, respectively. We see some evidence of positive ITE when
$\alpha$ is above $0.25$ while no evidence of any negative ITE even
when $\alpha = 0.5$.


%%% Local Variables:
%%% mode: latex
%%% TeX-master: t
%%% End:
